% Compile with: make   (runs make_report.py, then latexmk -pdf)
%
% Every quantity in the prose below is a macro from numbers.tex, and every
% table body is \input from tables/. Both are generated by make_report.py from
% the run's own output, so nothing here can drift from the data. If a number
% looks wrong, fix the generator, not this file.
\documentclass[11pt,a4paper]{article}

\usepackage[T1]{fontenc}
\usepackage[utf8]{inputenc}
\usepackage[margin=2.4cm]{geometry}
\usepackage{booktabs}
\usepackage{graphicx}
\usepackage[font=small,labelfont=bf]{caption}
\usepackage{microtype}
\usepackage{xcolor}
\usepackage{amsmath}
\usepackage[hidelinks]{hyperref}
\usepackage{titlesec}

\titlespacing*{\section}{0pt}{1.4ex plus .2ex}{0.7ex}
\titlespacing*{\subsection}{0pt}{1.1ex plus .2ex}{0.5ex}
\setlength{\parskip}{0.45em}
\setlength{\parindent}{0pt}
\renewcommand{\arraystretch}{1.12}

\input{numbers}

\title{\vspace{-1.6cm}Docking selectivity of \nMolecules{} generated molecules
against \nBacterial{} bacterial targets and human DHFR}
\author{M. Kikolski}
\date{\today}

\begin{document}
\maketitle
\vspace{-0.8cm}

\section{Summary}

\nInput{} molecules from the \texttt{exp\_c5\_admet\_mpo\_guided} set were
docked against \nReceptors{} receptors: \nBacterial{} bacterial targets
(\bacterialList) and human DHFR as a counter-screen.
\nScrubFailed{} molecules were lost during ligand preparation, leaving
\nMolecules{} molecules and \nDockings{} dockings, all of which completed.

The panel was not designed against these targets. It was produced by a
target-agnostic generative pipeline under an ADMET surrogate reward, and the
receptors were selected retrospectively, to ask whether the set is chemically
and geometrically compatible with validated antibacterial sites. Nothing below
is a prediction of activity, and no experimental measurement is involved at any
point.

Against the counter-screen the set is not selective. Comparing the two DHFR
orthologues directly --- the same fold with the same cofactor, both receptors
validated by redocking --- \textbf{\selDhfrN{} of \nMolecules{} molecules
score better against the bacterial enzyme}. The median difference is
\selDhfrMedian{}~kcal\,mol$^{-1}$ in favour of the human enzyme, and the
upper quartile is \selDhfrQHi{}, so the whole distribution lies on the same
side.

Widening the comparison to the best of any bacterial target,
\selAnyN{}~molecules (\selAnyPct\%) beat human DHFR at some target,
\selAnyHalf{} of them by at least \marginMin{}~kcal\,mol$^{-1}$ and
\selAnyOne{} by at least 1.0. These are cross-fold comparisons and so are
weaker evidence than the orthologue pair; they are also concentrated in a
single receptor (\S\ref{sec:nearmiss}).

The pattern is the same across all \nSeeds{} generator seeds
(\S\ref{sec:seeds}), so it does not come from one sampling run.

\section{Target selection}
\label{sec:targets}

This section is editorial rather than measured: it records why these
\nReceptors{} receptors were chosen, so the choice can be defended
independently of what the scores turned out to be. No target was defined
before generation, so the selection is retrospective by construction.

\subsection{What the panel allows}

Three properties of the panel constrain which sites can honestly be docked
into.

\emph{Nitrogen-rich, but not amine-guaranteed.} The generator filtered on
\emph{total} nitrogen count, which counts amide, nitrile and non-basic
aromatic nitrogen. It does not establish the ionisable, preferably primary
amine that the eNTRy rules identify as necessary for accumulation in
Gram-negative bacteria, alongside low globularity and few rotatable
bonds~\cite{richter2017}. Sites were therefore preferred where
\emph{aromatic and heterocyclic} nitrogen is itself the recognition element,
rather than sites that read a cation.

\emph{Small and synthesisable.} Sites whose validated ligands are large
lipophilic macrocycles binding shallow surfaces that refold on binding --- the
RNA polymerase switch region --- or charged phosphonate and nucleotide mimics
--- MurA--F, DXR/MEP --- are a poor match for the panel's property space.

\emph{Non-covalent, no metal centre.} Vina is a non-covalent empirical
function with no explicit electrostatic term, so covalent mechanisms
($\beta$-lactams at PBPs, fosfomycin at MurA, benzoxaboroles at tRNA
synthetases) and zinc-chelating hydroxamate chemistry (LpxC, peptide
deformylase) cannot be modelled honestly, whatever score they would
return. Table~\ref{tab:rejected} lists what this excluded.

\subsection{Structure selection}

Target choice and structure choice are separate decisions. Every entry in
Table~\ref{tab:structures} is \textbf{holo}, with a drug-like co-crystal
ligand in the pocket: apo structures present side chains and loops in a
conformation that is not binding-competent, and docking enrichment degrades
measurably against them~\cite{mcgovern2003}. Cofactors are retained rather
than stripped, because in both DHFR structures and in FabI the cofactor forms
part of the pocket wall --- trimethoprim binding to DHFR is NADPH-cooperative,
and triclosan stacks directly on the nicotinamide ring. Removing it would
leave a larger, emptier, physically wrong site. Redocking each entry's own
co-crystal ligand is the gate on all of this
(\S\ref{sec:methods}); it is also the only check that catches a cofactor or
ligand residue name mis-set in the configuration, which otherwise fails
silently.

\begin{table}[tb]
\centering
\caption{The structures docked, and the published source of each. 2W3A has no
primary publication; it is cited as a PDB deposition. Resolutions are as
reported by the RCSB entry.}
\label{tab:structures}
\small
\begin{tabular}{llcll}
\toprule
role in the panel & structure & res.\ / \AA & co-crystal ligand & source \\
\midrule
GyrB ATP site, Gram-positive & 3U2D \textit{S.\,aureus} & 1.85 &
pyrrolamide & \cite{eakin2012} \\
GyrB ATP site, Gram-negative & 4ZVI \textit{E.\,coli} & 2.20 &
\textit{N}-phenylpyrrolamide & \cite{zidar2015} \\
DHFR, bacterial & 7MYM \textit{E.\,coli} & 3.04 &
trimethoprim $+$ NADPH & \cite{krucinska2022} \\
DHFR, host counter-screen & 2W3A human & 1.50 &
trimethoprim $+$ NADPH & \cite{pdb2w3a} \\
FabI, Gram-positive pathogen & 3GR6 \textit{S.\,aureus} & 2.28 &
triclosan $+$ NADP & \cite{priyadarshi2010} \\
\bottomrule
\end{tabular}
\end{table}

\subsection{GyrB ATP site --- the broad-spectrum anchor}

The GyrB/ParE ATPase site is the most thoroughly worked unexploited
antibacterial site in the literature: novobiocin reached the clinic and was
withdrawn, and more than fifty years of pyrrolamide, aminobenzimidazole,
azaindole and pyridylurea programmes have followed without a registered
drug~\cite{bisacchi2015}. The domain is conserved across Gram-positives and
Gram-negatives and dual GyrB/ParE inhibition is an active strategy, which is
what makes it the broad-spectrum anchor here.

It is also the best \emph{chemical} match in the panel. Recognition runs
through a hydrogen bond to Asp73 (\textit{E.\,coli} numbering), frequently
bridged by a conserved structural water, donated by pyrrole NH and aromatic
nitrogen --- no cationic amine is required, which matters because the panel
cannot guarantee one. Both orthologues are included so that compatibility at
the fold can be separated from compatibility with one species' pocket.

The known pitfall is that several ligand series recruit that bridging water,
and omitting it deletes a real hydrogen-bond partner. This was tested rather
than assumed: 3U2D has four ordered waters within 3.5~\AA{} of the ligand, the
closest at 2.72~\AA, and with all of them excluded it still redocks at
0.53~\AA. 4ZVI has only a marginal 3.33~\AA{} contact. For these two entries
the water was not required.

\subsection{The DHFR orthologue pair --- the selectivity axis}

DHFR carries the only selectivity comparison available without a wet lab.
Trimethoprim discriminates the bacterial from the human enzyme by roughly
three to four orders of magnitude, an effect amplified by cooperative binding
with NADPH~\cite{wrobel2020}. That makes the orthologue pair a controlled
comparison in the sense docking can actually support: the same ligands, the
same fold, the same cofactor, and two receptors that differ in the pocket
rather than in everything at once. A score difference between unrelated folds
confounds the comparison with Vina's own calibration; a score difference
between two DHFRs does not.

The recognition motif suits the panel as well as GyrB does. The
2,4-diaminopyrimidine head group that engages the conserved Asp27 is
intrinsically nitrogen-rich, which is the closest correspondence in the whole
set between the generator's design heuristic and a validated pocket. This is
why the \selDhfrN{}/\nMolecules{} result in \S\ref{sec:result} is reported as
the headline rather than the cross-fold comparison: it is the test the panel
was best placed to pass.

\subsection{FabI --- the pathogen-relevant arm}

FabI, enoyl-ACP reductase, is the Gram-positive pathogen target of the set.
Triclosan is the tool compound, and afabicin (Debio 1450) has reached phase 2
trials against \textit{S.\,aureus}, including bone and joint infection
(NCT03723551) --- the most clinically advanced target here. Flat
nitrogen-containing aromatics stack against the NADP nicotinamide with a
hydrogen bond to Tyr156, which is compatible with the panel's heteroaromatic
content. Only ternary holo structures are usable, because the
substrate-binding loop is disordered otherwise.

Its weakness is stated rather than left for review. Whether type-II fatty acid
synthesis is a viable antibacterial target in Gram-positive pathogens is
contested: Brinster \textit{et al.}\ argued that these organisms bypass FAS-II
inhibition by scavenging host fatty acids~\cite{brinster2009}, and the
rebuttals --- including evidence that the result does not generalise to
\textit{S.\,aureus} --- leave the question
open~\cite{parsons2011}. This bears directly on the results below, because the
margins in \S\ref{sec:nearmiss} fall almost entirely on FabI.

\begin{table}[tb]
\centering
\caption{Targets considered and excluded, with the reason in each case. The
first three rows are excluded by the method, not by the biology.}
\label{tab:rejected}
\small
\begin{tabular}{p{0.24\linewidth}p{0.66\linewidth}}
\toprule
target & reason for exclusion \\
\midrule
PBP2a/PBP3 & $\beta$-lactam mechanism is covalent; non-covalent docking would
be misleading \\
MurA--MurF & large, polar, often charged nucleotide and UDP-sugar sites;
MurA inhibition by fosfomycin is covalent \\
LpxC, peptide deformylase & metalloenzymes whose potent inhibitors are
chelating hydroxamates; Vina has no explicit electrostatic term and mishandles
metal coordination \\
RNA polymerase switch region & large lipophilic macrocycles on a shallow
surface that locally refolds on binding \\
LptD/LptA, LolE & compelling biology, but identified phenotypically rather
than structurally; no clean high-resolution small-molecule pocket \\
tRNA synthetases & benzoxaborole mechanism is covalent and tRNA-dependent \\
DXR/MEP, UPPS, riboswitches & charged phosphonate sites, RNA targets, or thin
structural support \\
\bottomrule
\end{tabular}
\end{table}

\subsection{A possible extension: FabL}

One target was planned and not run. \textit{B.\,subtilis} FabL is not a FabI
variant: the two enzymes coexist in that organism, differ in cofactor (NADPH
against NADH) and share only around 24\% sequence identity, FabL being closer
to FabG and the wider SDR family while arriving at the same
Tyr-(Xaa)$_6$-Lys dyad~\cite{kim2011}. It is also the resistance determinant
--- \textit{fabI} knockouts are as triclosan-sensitive as wild type, while
\textit{fabL} knockouts are some 250-fold more
sensitive~\cite{heath2000}. Two things would make it worth adding. It is
unusually tractable for docking, because triclosan inhibits FabI through a
slow-binding, effectively mechanism-based ternary complex, whereas FabL is
inhibited reversibly, and ordinary reversible occupancy is what a scoring
function can represent. And FabI against FabL would give a cross-isoform
compatibility profile --- structurally the same question as the DHFR pair
above, but about antibacterial spectrum rather than host selectivity. The
highest-quality structure in the whole candidate set, 3OID at 1.80~\AA,
is available for it. It is an extension for comparison, not a gap in the
present result: nothing reported here depends on it.

\section{Result}
\label{sec:result}

\subsection{Selectivity against human DHFR}

Table~\ref{tab:selectivity} gives both comparisons. Figure~\ref{fig:sel}
shows the orthologue difference per molecule: the distribution lies almost
entirely on the side favouring the human enzyme, and its upper tail does not
cross zero. Figure~\ref{fig:rec} shows the same effect as score
distributions, with human DHFR shifted towards stronger binding than all
\nBacterial{} bacterial targets.

\begin{table}[tb]
\centering
\caption{Molecules favouring a bacterial target over human DHFR. Positive
values mean the bacterial target scores more strongly. The first row is the
comparison the data supports directly; the second spans different folds.}
\label{tab:selectivity}
\small
\input{tables/selectivity}
\end{table}

\begin{figure}[tb]
\centering
\includegraphics[width=0.86\linewidth]{figures/selectivity}
\caption{Per-molecule difference between the human and \textit{E.\,coli} DHFR
scores. Values right of the dashed line would favour the bacterial enzyme.}
\label{fig:sel}
\end{figure}

\begin{figure}[tb]
\centering
\includegraphics[width=0.86\linewidth]{figures/receptors}
\caption{Best Vina score per molecule at each receptor, stronger binding
upwards. Human DHFR (right, red) is shifted towards stronger binding than
every bacterial target in the panel.}
\label{fig:rec}
\end{figure}

\subsection{Ligand efficiency}

The strongest ligand efficiency in the panel is \leBest{} at
\leBestTarget; among the bacterial targets the best is \leBestBacterial{} at
\leBestBacterialTarget. Both come from the smallest molecule in the set.
Taking \leFloor{} as a conventional floor~\cite{hopkins2004} for a hit
worth progressing, no
molecule combines that efficiency with a preference for a bacterial target.

Vina scores are often inflated for larger ligands, which would distort a
ranking of this kind. That effect is not present here: the correlation between
score and heavy-atom count runs from \sizeCorrLo{} to \sizeCorrHi{} across the
\nReceptors{} receptors, i.e.\ marginally positive. The heavy-atom range is
narrow (median \heavyMedian, from \heavyMin{} to \heavyMax), which is the
likely reason. Raw scores and efficiencies therefore rank this particular set
consistently.

\begin{figure}[tb]
\centering
\includegraphics[width=0.86\linewidth]{figures/efficiency}
\caption{Ligand efficiency at each molecule's best bacterial target against
molecular size. Circled points are the \nNearMiss{} molecules that beat human
DHFR by at least \marginMin{}~kcal\,mol$^{-1}$.}
\end{figure}

\subsection{Consistency across generator seeds}
\label{sec:seeds}

Table~\ref{tab:seeds} splits the margin by the generator seed recorded in each
molecule's identifier. The medians differ by at most
\seedGap{}~kcal\,mol$^{-1}$, and a permutation test on that widest pair gives
$p=\seedPval$, so the seeds are not distinguishable on this measure.

\begin{table}[tb]
\centering
\caption{Selectivity margin against human DHFR, by generator seed. A
permutation test on the widest pair of medians (seed \seedHi{} versus
\seedLo, gap \seedGap{}~kcal\,mol$^{-1}$) gives $p=\seedPval$ over
\seedShuffles{} shuffles.}
\label{tab:seeds}
\small
\input{tables/seeds}
\end{table}

\subsection{Molecules with any selectivity margin}
\label{sec:nearmiss}

Table~\ref{tab:nearmiss} lists every molecule beating human DHFR by at least
\marginMin{}~kcal\,mol$^{-1}$, at \nearMissTargets. None reaches the
\leFloor{} efficiency floor. Section~\ref{sec:scope} notes why these
particular targets carry less weight than the orthologue comparison.

\begin{table}[tb]
\centering
\caption{All \nNearMiss{} molecules beating human DHFR by at least
\marginMin{}~kcal\,mol$^{-1}$, at \nearMissTargets. None reaches the
\leFloor{} efficiency floor. The best of them by efficiency is
\bestLECandidate{} (\bestLECandidateTarget,
\bestLECandidateScore~kcal\,mol$^{-1}$, \bestLECandidateHeavy{} heavy atoms,
LE \bestLECandidateLE, margin \bestLECandidateMargin). Identifiers are
abbreviated to seed/serial; full names are in the result CSVs.}
\label{tab:nearmiss}
\small
\input{tables/nearmiss}
\end{table}

\section{Methods and validation}
\label{sec:methods}

Ligands were enumerated to protonation states and 3D conformers with molscrub
at pH 7.4 and typed with Meeko. Receptors were taken from the RCSB entries in
Table~\ref{tab:receptors}, reduced to a single chain, protonated with PDBFixer
at pH 7.4 and typed with Meeko. Cofactors (NADP$^+$ in the two DHFR structures
and in FabI) were prepared separately and appended as rigid atoms, because
Meeko carries no residue template for them. Docking used AutoDock Vina 1.2~\cite{vina2021}
through its Python API at exhaustiveness 32, nine retained poses, seed 2137,
with boxes of 20--22~\AA{} centred on each entry's co-crystal ligand.

Each receptor was validated by redocking its own co-crystal ligand, with bond
orders taken from the chemical component's idealised structure rather than
perceived from coordinates. The gate is recovery of the crystal pose within
\redockThreshold~\AA{} among the nine retained poses, which all
\nRedockTotal{} receptors pass. The rank at which that pose appears is
reported separately, as a per-receptor confidence weight rather than a
pass/fail criterion: \nRedockPass{} of \nRedockTotal{} also place it first.

\begin{table}[tb]
\centering
\caption{Panel, score distribution and redocking validation. RMSD values are
in \AA; ``best'' is the closest of the nine retained poses and ``rank'' its
position in Vina's ordering.}
\label{tab:receptors}
\small
\input{tables/receptors}
\end{table}

\section{Scope of the conclusion}
\label{sec:scope}

\begin{itemize}\setlength{\itemsep}{0.2em}
\item The orthologue comparison (DHFR \textit{E.\,coli} against human DHFR)
is the one this data supports directly: identical ligands, the same fold, the
same cofactor, and both receptors passing the redocking gate.
\item Comparisons between different folds --- FabI or GyrB against human DHFR
--- are weaker. Vina scores are calibrated for ranking ligands at one site,
not for comparing absolute affinities between unrelated sites.
\item Two receptors, FabI \textit{S.\,aureus} and GyrB \textit{E.\,coli}, do
not place the crystal pose first, so their scores carry less weight than the
other three. The margins in \S\ref{sec:nearmiss} fall almost entirely on
FabI.
\item 7MYM, which carries the headline comparison, is the lowest-resolution
entry in the panel at 3.04~\AA, against 1.50~\AA{} for its human counterpart.
It recovers its own co-crystal pose at 0.89~\AA{} and ranks it first, so the
site is modelled adequately, but a higher-resolution \textit{E.\,coli}
DHFR:trimethoprim:NADPH entry would be the first thing to substitute on a
re-run.
\item FabL \textit{B.\,subtilis} (3OID) was planned and not docked
(\S\ref{sec:targets}). Nothing above depends on it; it would add a
cross-isoform comparison, not close a gap.
\item No actives-versus-decoys enrichment test (ROC AUC, EF1\%, BEDROC) has
been run, so by the project's own selection criteria every target here remains
provisional.
\item Whether type-II fatty acid synthesis is a viable target in
Gram-positive pathogens is contested in the literature. That bears directly on
the FabI results, which carry most of the margins reported above.
\item \nScrubFailed{} of \nInput{} molecules failed ligand preparation and are
absent from every figure here.
\item Vina scores are a ranking device, not predicted binding free
energies: correlation with measured affinity is weak across docking programs
and scoring functions~\cite{warren2006}, and no experimental measurement is
involved at any point.
\end{itemize}

\begin{thebibliography}{99}
\small
\bibitem{eakin2012} Eakin A.E. \textit{et al.} Pyrrolamide DNA gyrase
inhibitors: fragment-based NMR screening to identify antibacterial agents.
\textit{Antimicrob.\ Agents Chemother.} \textbf{56} (2012) 1240. PDB 3U2D.

\bibitem{zidar2015} Zidar N. \textit{et al.}
\textit{N}-Phenyl-4,5-dibromopyrrolamides and \textit{N}-phenylindolamides as
ATP-competitive DNA gyrase B inhibitors. \textit{J.\ Med.\ Chem.}
\textbf{58} (2015) 6179. PDB 4ZVI.

\bibitem{krucinska2022} Krucinska J. \textit{et al.} Structure-guided
functional studies of plasmid-encoded dihydrofolate reductases reveal a common
mechanism of trimethoprim resistance in Gram-negative pathogens.
\textit{Commun.\ Biol.} \textbf{5} (2022) 459. PDB 7MYM.

\bibitem{pdb2w3a} PDB 2W3A, human dihydrofolate reductase in complex with
NADPH and trimethoprim, 1.50~\AA. Deposited 2008; no primary publication.

\bibitem{priyadarshi2010} Priyadarshi A., Kim E.E., Hwang K.Y. Structural
insights into \textit{Staphylococcus aureus} enoyl-ACP reductase (FabI) in
complex with NADP and triclosan. \textit{Proteins} \textbf{78} (2010) 480.
PDB 3GR6.

\bibitem{bisacchi2015} Bisacchi G.S., Manchester J.I. A new-class
antibacterial --- almost. Lessons in drug discovery and development: a
critical analysis of more than 50 years of effort toward ATPase inhibitors of
DNA gyrase and topoisomerase IV. \textit{ACS Infect.\ Dis.} \textbf{1} (2015)
4.

\bibitem{wrobel2020} Wróbel A., Arciszewska K., Maliszewski D., Drozdowska D.
Trimethoprim and other nonclassical antifolates: an excellent template for
searching modifications of dihydrofolate reductase enzyme inhibitors.
\textit{J.\ Antibiot.} \textbf{73} (2020) 5. doi:10.1038/s41429-019-0240-6.

\bibitem{brinster2009} Brinster S. \textit{et al.} Type II fatty acid
synthesis is not a suitable antibacterial target for Gram-positive pathogens.
\textit{Nature} \textbf{458} (2009) 83.

\bibitem{parsons2011} Parsons J.B., Rock C.O. Is bacterial fatty acid
synthesis a valid target for antibacterial drug discovery?
\textit{Curr.\ Opin.\ Microbiol.} \textbf{14} (2011) 544.

\bibitem{kim2011} Kim K.-H. \textit{et al.} Crystal structures of enoyl-ACP
reductases I (FabI) and III (FabL) from \textit{Bacillus subtilis}.
\textit{J.\ Mol.\ Biol.} (2011). doi:10.1016/j.jmb.2010.12.003. PDB 3OID.

\bibitem{heath2000} Heath R.J. \textit{et al.} The
enoyl-[acyl-carrier-protein] reductases FabI and FabL from \textit{Bacillus
subtilis}. \textit{J.\ Biol.\ Chem.} \textbf{275} (2000) 40128.

\bibitem{richter2017} Richter M.F. \textit{et al.} Predictive compound
accumulation rules yield a broad-spectrum antibiotic. \textit{Nature}
\textbf{545} (2017) 299.

\bibitem{mcgovern2003} McGovern S.L., Shoichet B.K. Information decay in
molecular docking screens against holo, apo, and modeled conformations of
enzymes. \textit{J.\ Med.\ Chem.} \textbf{46} (2003) 2895.

\bibitem{vina2021} Eberhardt J., Santos-Martins D., Tillack A.F., Forli S.
AutoDock Vina 1.2.0: new docking methods, expanded force field, and Python
bindings. \textit{J.\ Chem.\ Inf.\ Model.} \textbf{61} (2021) 3891.

\bibitem{hopkins2004} Hopkins A.L., Groom C.R., Alex A. Ligand efficiency: a
useful metric for lead selection. \textit{Drug Discov.\ Today} \textbf{9}
(2004) 430.

\bibitem{warren2006} Warren G.L. \textit{et al.} A critical assessment of
docking programs and scoring functions. \textit{J.\ Med.\ Chem.} \textbf{49}
(2006) 5912.
\end{thebibliography}

\appendix
\section{Provenance}

Pipeline commit \texttt{9770a4e}, recorded in the job log and in every stage
manifest alongside package versions, resolved tool paths and SHA-256 checksums
of the inputs. The run was SLURM job 1994, completed 2026-10-18, taking 134~h
on one 32-core node at a mean 104.8~min per docking. Figures and tables in
this document were generated by \texttt{report/make\_report.py} from
\texttt{analyze.py} output, with that script's default thresholds
(LE $\geq$ \leFloor, margin $\geq$ \marginMin{}~kcal\,mol$^{-1}$).

\end{document}
